\subsection{An explicit finite-field determination}

The algebraically closed construction has a finite-field counterpart.
Here a finite search replaces eigenvalue ordering, and the multiplicity
forms contribute explicit two-torsion. We describe all those searches,
including the parameter set, so that an unknown representation-isomorphism
index is not left in the answer.

\begin{theorem}[Explicit canonical indices and operations over every finite field]
\label{C:Om-finite-field-algorithm}
\status{Proved}
Let $k=\mathbb F_q$ with $q$ odd, let $m\geq0$, and fix a split
nonsingular symmetric matrix $J$ defining $O_m$. Fix an ordering of
the finite set $k$. The finite procedure below produces a repetition-free
index set $\mathscr J$ for every indecomposable rational $O_m$-module,
and computes duality and its multiplicity-form data. For each
dimension $n$ and minimal coefficient degree $d$ the list is finite.

Let $U_i$ be its indexed modules. For each self-dual $i$, choose the first
invariant nonsingular $\varepsilon_i$-symmetric form in the finite
enumeration of form matrices, searching the symmetric case first.
Write $D_i=\End_{O_m}(U_i)/\operatorname{rad}$ and let $\sigma_i$ be
its adjoint involution. These are explicitly computed finite fields
with involution. Partition the indices into
\[
 \begin{array}{c|l}
 \mathscr J_{\mathrm s}&U_i\simeq U_i^*,\quad
       \sigma_i=1,\quad\varepsilon_i=1,\\
 \mathscr J_{\mathrm a}&U_i\simeq U_i^*,\quad
       \sigma_i=1,\quad\varepsilon_i=-1,\\
 \mathscr J_{\mathrm u}&U_i\simeq U_i^*,\quad\sigma_i\ne1,\\
 \mathscr J_{\mathrm h}&\text{the smaller index in each non-self-dual pair.}
 \end{array}
\]
Then there is the explicit additive decomposition
\[
 WG_{\mathrm{sp}}(BO_m)=
 \bigoplus_{i\in\mathscr J_{\mathrm s}}
       (\mathbb Z e_i\oplus(\mathbb Z/2)t_i)
 \ \oplus\!
 \bigoplus_{i\in\mathscr J_{\mathrm a}\cup
                    \mathscr J_{\mathrm u}\cup\mathscr J_{\mathrm h}}
       \mathbb Z e_i.                                      \tag{C.3.F1}
\]
All generators have the concrete representatives specified in the proof.
Their products, and every exterior-power coefficient of any finite
integral combination, are determined by terminating finite procedures.
Thus (C.3.F1), with those procedures, determines the entire split ring
and its exterior operations over $k$.
\end{theorem}
\begin{proof}
\emph{An index involving only finite matrix enumeration.}
Construct $H=k[O_m]$ and its nested standard-monomial coalgebras $C_d$
as in the proof of Theorem~\ref{C:Om-canonical-templates}. That construction
uses only field arithmetic and Gr\"obner reduction, and works over $k$
without algebraic closure. Put $r=\dim_k C_d$ and compute its dual-algebra
structure constants. Enumerate all $q^{rn^2}$ tuples in $M_n(k)^r$.
Retain those satisfying the explicit dual-algebra multiplication and unit
equations, and whose coefficients outside $C_{d-1}$ do not all vanish
when $d>0$. For each retained tuple, enumerate the at most $q^{n^2}$
matrices in $M_n(k)$ and discard the tuple if its centralizer contains
an idempotent other than $0,I$. Finally retain only the lexicographically
least simultaneous conjugate under the finite group $\mathrm{GL}_n(k)$.
All tests are polynomial matrix equalities or determinants in the finite
field. The union of these finite lists over $(n,d)$ is $\mathscr J$.

The coaction equivalence, idempotent criterion and invariance of minimal
coefficient degree proved in Theorem~\ref{C:Om-canonical-templates}
show exhaustiveness and absence of repetition. The least representative
is determined by a search over a specified finite set, not by a further
isomorphism quotient. The antipode formula produces the dual tuple;
the same finite conjugation search gives its canonical index.

\emph{The finite multiplicity algebras and generators.}
Compute $E_i=\End_{O_m}(U_i)$ by the coefficient-commutator linear system.
It is a finite-dimensional local algebra, by the indecomposability and
Fitting argument in Proposition~\ref{C:field_WG_structure}. Its nonunits
are exactly its radical and are found by checking determinants of its
finite set of elements. Form the quotient multiplication table $D_i$.
This is a finite division ring, hence a field by Wedderburn's theorem
in the opened source \cite[Theorem III.4.4, p.29]{C-ArtinFiniteDivision}.
For a self-dual module, Proposition~\ref{C:field_WG_structure} guarantees
that one of the two finite searches for nonsingular symmetric or
alternating invariant matrices succeeds. Its adjoint on $E_i$ descends
to the explicitly computed involution $\sigma_i$ on $D_i$.

We now apply the hermitian-category decomposition of
Proposition~\ref{C:field_WG_structure} with these particular indices.
The finite-field form classifications in
\cite[Lemma 3.1(a),(b)]{C-SergeichukFiniteForms} give
\[
 \begin{array}{c|c}
 (D,\sigma,\varepsilon)&GW^\varepsilon(D,\sigma)\\ \hline
 (D,1,1)&\mathbb Z\oplus\mathbb Z/2\\
 (D,1,-1)&\mathbb Z\\
 (D,\sigma\ne1,\pm1)&\mathbb Z.
 \end{array}                                               \tag{C.3.F2}
\]
We spell out the deduction to fix the split, rather than Witt, convention.
In the first case a nonsingular form is classified by its dimension and
the square class of its determinant. If $c$ is the first nonsquare in
$D$, the group completion is generated by $\langle1\rangle$ and
$\langle c\rangle-\langle1\rangle$, the latter of exact order two.
In the second case the dimension is even and determines the alternating
form; a symplectic plane is one generator. For nonidentity $\sigma$,
the hermitian dimension determines the form. A nonzero $a$ with
$\sigma(a)=-a$ changes skew-hermitian to hermitian forms when needed.
Such $a$ exists by taking $b-\sigma(b)$ with $b$ not fixed by $\sigma$.
Thus this case also has one generator, with multiplicity dimension one.
Non-self-dual pairs give the hyperbolic integer summands. These are
exactly the summands of (C.3.F1).

Every generator can be represented by finite linear algebra. Let $q_i$
be the chosen form and $*$ its adjoint on $E_i$. For
$i\in\mathscr J_{\mathrm s}$ take $e_i=[q_i]$. Lift the first nonsquare
$c\in D_i$ to $b\in E_i$ and put $a=(b+b^*)/2$. It is self-adjoint
and invertible, since its reduction is $c$. The invariant symmetric
form $q_i a$ defines $s_i$, and put $t_i=s_i-e_i$.
For $i\in\mathscr J_{\mathrm a}\cup\mathscr J_{\mathrm h}$ take
$e_i=H(U_i)$. For $i\in\mathscr J_{\mathrm u}$ choose the first nonzero
$c\in D_i$ with $\sigma_i(c)=\varepsilon_i c$, lift it to $b$, and put
$a=(b+\varepsilon_i b^*)/2$. Then $q_i a$ is symmetric and invertible;
its class is $e_i$. The residue-form classification and the nilpotent
radical lifting theorem already checked in
Proposition~\ref{C:field_WG_structure} prove that these representatives
have precisely the claimed generators and relations.

\emph{Operations and termination.}
The tensor coaction and exterior coactions are the products and minors
of their coefficient matrices. Gr\"obner reduction places them in
$C_{d+e}$ and $C_{ad}$ for a tensor and an $a$th exterior power,
respectively. Their Gram matrices are tensor products and the induced
determinant forms. No factorial is inverted.

Given any resulting actual representation of dimension $N$ and coefficient
degree at most $D$, split nontrivial commuting idempotents by the finite
matrix search until every summand is indecomposable. Its dimensions
strictly decrease at every split. Compute each summand's minimal degree
and canonical index. To determine the form, the underlying multiplicities
leave only finitely many choices: one determinant bit for each occurring
index in $\mathscr J_{\mathrm s}$, and the forced multiplicities in the
other three types. Construct the corresponding orthogonal sums of the
displayed generators and test equivariant isometry by the finite search
through $\mathrm{GL}_N(k)$. The hermitian-category theorem and (C.3.F2)
ensure that exactly one of these choices succeeds. Its coordinates give
the product or exterior-power coefficients in (C.3.F1).

For example, multiplicity $b>0$ and determinant bit $c\in\{0,1\}$
in symmetric type give the actual form consisting of $b-c$ copies of
$e_i$ and $c$ copies of $s_i$, with group coordinate $be_i+ct_i$.
No determinant bit is assigned to a zero multiplicity. Alternating type
has even underlying multiplicity and uses half as many hyperbolic
generators. The two members of a non-self-dual pair have equal
multiplicities. These facts also make all search bounds explicit.

Bilinearity treats virtual products, and
$\lambda_t(x-y)=\lambda_t(x)\lambda_t(y)^{-1}$ treats virtual exterior
powers to any requested finite degree. Each step consists of a finite
enumeration or polynomial reduction on specified finite matrices. This
proves the claimed determination, without a complexity bound.
\end{proof}

\begin{proposition}[The finite-field form parameter on the characteristic-three example]
\label{C:Om-finite-field-form-example}
\status{Proved}
Over $\mathbb F_3$, let $S=\operatorname{Sym}^2(\mathbb F_3^3)$ with
its rational $O_3$-action and invariant form $Q(A,B)=\operatorname{tr}(AB)$.
The forms $Q$ and $2Q$ are not equivariantly isometric. Their difference
is a nonzero element of exact order two in $WG_{\mathrm{sp}}(BO_3)$,
although the underlying representations agree. The isotypic contribution
of $S$ is $\mathbb Z\oplus\mathbb Z/2$.
\end{proposition}
\begin{proof}
The coefficient computation in
Proposition~\ref{C:Om-canonical-example} is defined over $\mathbb F_3$.
It gives $\End_{O_3}(S)=\mathbb F_3[I,K]$ with $K^2=0$, and the
nontrivial matrix $K=q\operatorname{tr}$ is self-adjoint for $Q$.
Consequently every equivariant automorphism is $u=aI+bK$ with $a\ne0$,
and its residual norm is $a^2=1$. It cannot carry $Q$ to $2Q$.
More generally every invariant nonsingular symmetric form has matrix
$Q(cI+dK)$, where $c\ne0$, and its isometry class is determined by
$c\in\mathbb F_3^\times/(\mathbb F_3^\times)^2$. The nilpotent
coefficient is removed by the $2abK$ term in $u^*u$.
Thus the multiplicity involution is the identity on $D=\mathbb F_3$,
and (C.3.F2) gives the stated group, including nonzero two-torsion.
Concretely $\left(\begin{smallmatrix}1&1\\-1&1\end{smallmatrix}\right)$
has transpose-product $2I_2$ over $\mathbb F_3$; its tensor with $I_S$
gives $Q\perp Q\simeq2Q\perp2Q$. The exact residue determinant
invariant prevents the difference from vanishing after stabilization.
The certificate \texttt{code/C\_finish\_finite\_field\_forms.py}
enumerates all $6$ equivariant units and all $6$ nonsingular invariant
symmetric forms in the $9$-element endomorphism algebra. Its $36$
action comparisons give exactly two orbits of size three, and it
checks the displayed double-sum congruence in dimension $12$.
The full endomorphism-algebra calculation is independently included
in the earlier coordinate-coalgebra certificate.
\end{proof}

This finite-field theorem and the algebraically closed theorem are
separate parts of C.3. Neither supplies a classification over every
infinite nonclosed field or over every nonfield base with two invertible.
